\section{客观违法阻却事由 - 紧急避险}
\begin{enumerate}
  \item 基本同正当防卫。正当防卫针对不法侵害人本人，紧急避险对第三方。
  \item 成立条件。
        \begin{enumerate}
          \item 起因条件。一份\zd{单纯的危险}。
                \begin{enumerate}
                  \item \al{野狗恶犬追咬（单纯的风险），躲到别人家}。
                  \item \al{别人拿刀砍（不法侵害）追你，躲到别人家}。
                \end{enumerate}
          \item 时间条件。基本同正当防卫。
          \item 意思条件。
                \begin{enumerate}
                  \item 避险意图/动机。不要求具有。
                  \item 避险认识。不要说和必要说。
                        \begin{enumerate}
                          \item 假想避险。基本同假象防卫。
                          \item \zd{偶然避险}。\al{狗蛋为了泄愤砸车，孩子在车里差点被闷死}。
                        \end{enumerate}
                \end{enumerate}
          \item \zd{不得已条件/补充性手条件}。除了 A 类手段，没有其他合理的 B 类手段可以选择，\zd{不看主观}。
                \begin{enumerate}
                  \item 山穷水尽。\al{甲房屋失火，\scx{有乙丙丁鱼塘，甲就是看乙不顺眼，故意}抽乙鱼塘的水灭火}。
                  \item 受强制的紧急避险。\al{检察官被抓强迫强奸}。
                \end{enumerate}
          \item 限度条件。\zd{保护的法益 $\geq$ 损害的法益}。
        \end{enumerate}
\end{enumerate}
